\pg{351} \emph{acío} puro, não encerra em si \textbf{mesmo} nenhuma potencialidade. Ora o ente, que não encerra em si nenhuma potencialidade, é absolutamente \emph{simples}; porque é só acto, e o acto não tem paries, em que possa dividir-se; Logo o ente; que é \emph{acto puro}, é absolutamente \emph{simples}.

b) \emph{O ente}, \emph{que} é absolutamente, simples, é \emph{acto} puro. --- O ente, que é absolutamente \emph{simples}, não consta de partes, e por isso não é composto de potencia e de acto. Ora o ente, que não composto de potencia e de acto, é \emph{actio puro}. Logo o ente, que absolutamente \emph{simples}, \emph{é acto puro} \footnote{O acto puro é absolutamente \emph{simples}. O acto é perfeição, e o acto puro é só perfeição, é perfeição subsisteote: e por isso exclue toda 2a potencialidade, não pode ser composto de potencia e de acto, é absolutamente \textbar{} simples. Por quanto, como diz S. Thomaz, em todo o composto deve encontrar-se o acto e a potencia. Com effeito, muitas coisas não podem tornar-se simplesmente um: só coisa, se não eniquanto uma seja acto e outra potencia; visto que, se fossem todas em acto, haveria um aggregado de coisas, mas nunes uma só coisa composta. Ora o acto puro, excluindo toda e qualquer potencialidade, é absolutamente simples. --- Não é menos evidente que o ente absolutamente \emph{simples é acto} puro. Por quanto, o ente absolutamente \emph{simples}. que subsiste perfeitamente em si mesmo, ou é só potencia, ou é só telo. Ora € absurdo dizer que é só potencia; pois. como vimos, a potencia, denotando limitação e imperfeição, não \emph{pode} subsistir em si mesma. Logo deve ser só acto, deve ser acto puro (\emph{Cf}. \emph{C}. \emph{Gent}., 1, 18).}.

75. O \textbf{ente} composto não pode ser acto \textbf{puro}. --- Ente \emph{composto} é o que tem partes. Ora o que tem partes não pode ser

\emph{X acto puro}; porque --- ou uma das partes será \emph{acto} em relação á outra,---- ou todas as partes \textbf{serão} como que \emph{em potencia} em relação ao todo. Logo o \textbf{ente} \emph{composto} não pode ser \emph{acto puro} \footnote{Diz S. Thomaz: «In omni composilo oportet esse potentiam et acinm; \emph{quia vel} ana pars est actus respectu alterius, vel saltem omnes partes. sunt in potentia respeetu Lotius» (\emph{Sum}, \emph{Th}., \emph{p}. Y, q. 3, a. 7).}. =

76. O ente, que é absolutamente simples, é necessario e infinito; o ente, que é necessario e infinito, é absohitamente simples. :

a) \emph{O ente}, \emph{que é absolutamente simples}, \emph{é necessario} e \emph{infinito}. --- O ente, que é absolutamente \emph{simples}, é acto puro, é o.proprio ser subsistente. Ora o acto puro, o proprio ser subsistente, necessario e \emph{infinito}. Logo o ente que é absolutamente \emph{simples}, \textbar{} 6 \emph{necessario} e \emph{infinito}. '

b) \emph{O ente}, \emph{que é necessario e infinito}, \emph{é absolutamente simples}. --- O ente, que é \emph{necessario} e \emph{infinito}, possue a existencia 